\chapter*{Abstract}
\addcontentsline{toc}{chapter}{Abstract}

{\setstretch{1.2}
Canny Vision is a company specializing in computer vision solutions, whose
project teams require an efficient means of tracking tasks, managing goals, and resolving
technical issues across multiple concurrent projects. Existing commercial platforms such as
Jira or ClickUp were ruled out on two grounds: hosting sensitive project data on third-party
servers raises data-sovereignty and security concerns incompatible with the company's
zero-trust posture, and no off-the-shelf tool can be fully tailored to the company's specific
workflows and internal constraints.

This report presents \textbf{Canny TMS}, a custom-built Ticket Management System designed
and developed as an end-to-end internal solution for Canny Vision. The system integrates
seamlessly with the company's pre-existing Identity and Access Management infrastructure,
based on OAuth~2.1 with PKCE, ensuring that all authentication and authorization flows
remain under full organizational control. The frontend is implemented in pure vanilla
JavaScript using the company's in-house MVP rendering engine, with no reliance on
third-party libraries. A fine-grained Attribute-Based Access Control model, conformant
with the NIST SP~800-162 standard, enforces workspace-level authorization policies and
governs every protected operation within the system.

Canny TMS is fully operational and currently in production use across multiple teams at
Canny Vision. It supports the complete lifecycle of epics and tickets --- from creation
and assignment to status tracking and threaded commenting --- while maintaining strict
data isolation within a self-hosted deployment.

The system demonstrates how security-first design constraints and a zero-external-dependency
policy can be reconciled with a productive and user-friendly collaboration tool, and lays
the groundwork for future enhancements such as a Policy Administration Point and advanced
analytics dashboards.

\bigskip
\noindent\textbf{Keywords:} ticket management, zero-trust architecture, OAuth~2.1 with PKCE,
attribute-based access control, vanilla JavaScript.
}
